\subsection{Esp32}  

Συμφώνα με την ίδια την εταιρία που παράγει τους esp32,η Espressif,οι esp32 είναι συστήμα σε ένα ολοκληρωμένο σύστημα (System on Chip,SoC)με 
Το Esp32 αποτελέι ένας  πολύ δυνατός και γενικού τύπου μικροελλεκτής,που αναφέρεται σαν "συστήμα σε ένα ολοκληρωμένο κύκλωμα" ), που 

ESP32 is a single 2.4 GHz Wi-Fi-and-Bluetooth combo chip designed with the TSMC low-power 40 nm
technology.



ESP32 Functional Block Diagram

1.Preface
ESP-32S is a powerful, generic WiFi-BT-BLE MCU module that targets a wide variety of applications
ranging from low power sensor networks to the most demanding tasks such as voice encoding, music
streaming and MP3 decoding.
At the core of this module is the ESP32 chip, which is designed to be scalable and adaptive. There are 2
CPU cores that can be individually controlled or powered, and the clock frequency is adjustable from 80
MHz to 240 MHz. The user may also power off the CPU and make use of the low power coprocessor to
constantly monitor the peripherals for changes or crossing of thresholds. ESP32 integrates a rich set of
peripherals, ranging from capacitive touch sensors, Hall sensors, low noise sense amplifiers, SD card
interface, Ethernet, high speed SDIO/SPI, UART, I2S and I2C.
The integration of Bluetooth, Bluetooth LE and Wi-Fi ensures that a wide range of applications can be
targeted,and that it is future proof: using Wi-Fi allows a large physical range and direct connection to
the internet through a Wi-Fi router, while using Bluetooth allows the user to conveniently connect to
the phone or broadcast low energy beacons for its detection. The sleep current of the ESP32 chip is less
than 5 μA, making it suitable for battery powered and wearable electronics applications. ESP-WROOM-
32 supports data rates up to 150 Mbps, and 22 dBm output power at the PA to ensure the widest
physical range. As such the chip does offer industry leading specifications and the best optimized
performance for electronic integration, range and power consumption, and connectivity.
The operating system chosen for ESP32 is freeRTOS with LWIP; TLS 1.2 with hardware acceleration is
built in as well. Secure (encrypted) over the air (OTA) upgrade is also supported, so that developers can
continually upgrade their products even after their release. The software releases are covered under
the ESP32 bug bounty program nd any bugs can be reported to bugbounty@espressif.com. As the SDK
of ESP-32S or ESP32 is opensource, the user can build his own platforms and operating systems. For
more in-depth discussion of this, the developer can contact john.lee@espressif.com.


ESP32 is powerful SoC (System on Chip) microcontroller with integrated Wi-Fi 802.11 b/g/n, dual mode Bluetooth version 4.2 and variety of peripherals. 
It is an advanced successor of the 8266 chip primarily in the implementation of two cores clocked in different version up to 240 MHz.
 Compared to its predecessor, except these features, it also extends the number of GPIO pins from 17 to 36, the number of PWM channels per 16 and is equipped with 4MB 
 of flash memory.

The ESP32 chip has been developed by the Espressif Systems company, which currently offers several ESP32 versions of the SoC in the form of ESP32 Developer Kit, 
the ESP32 Wrover Kit, which also includes an SD card and 3.2” LCD display and last but not least the ESP32 Azure IoT kit with USB Bridge and other built-in sensors.
 In addition to Espressif Systems, other producers are devoted to these chips - SparkFun with ESP32 Thing DB, WeMoS with its TTGO, D1, Lolin32 and Lolin D32, 
 Adafruid (with Huzzah32), DF Robot (ESP32 FireBeeatle) and many other manufacturers sometimes offer good and sometimes bad clones.

ESP32 includes two core (Xtensa LX6 processor made with 40 nm technology). CPU cores can be individually controlled. 
There is 520 KB of on-chip SRAM for data and instructions available. Some SoC modules such as ESP32-Wrover features 4 MB of external SPI flash and an additional 8 MB of
 SPI PSRAM (Pseudo static RAM). We have the possibility to use SPI, I2S, I2C, CAN, UART, Ethernet MAC, and IR in various quantities, depending on the type of the board. |
 Standard equipment also includes Hall Effect sensor, temperature sensor and touch sensor, other built-in sensors are implemented in Azure IoT and Developer kit. 
 SoC also provides Cryptographic hardware acceleration: AES, SHA-2, RSA, and Elliptic Curve Cryptography (ECC) and random number generator (RNG). 
 Figure 1 shows the ESP32 microcontroller variants with which we are working at the Department of Control System and Instrumentation.

ESP32 boards are produced in prototype designs that can be used in smart home applications, automation, wearables, audio applications,
 cloud-based IoT applications, and more. It is possible to choose a specific development kit or to design a custom embedded system built on the ESP32 microcontroller. 




Features
Wi-Fi
• 802.11b/g/n
• 802.11n (2.4 GHz), up to 150 Mbps
• WMM
• TX/RX A-MPDU, RX A-MSDU
• Immediate Block ACK
• Defragmentation
• Automatic Beacon monitoring (hardware TSF)
• Four virtual Wi-Fi interfaces
• Simultaneous support for Infrastructure Station, SoftAP, and Promiscuous modes
Note that when ESP32 is in Station mode, performing a scan, the SoftAP channel will be changed.
• Antenna diversity
Bluetooth®
• Compliant with Bluetooth v4.2 BR/EDR and Bluetooth LE specifications
• Class-1, class-2 and class-3 transmitter without external power amplifier
• Enhanced Power Control
• +9 dBm transmitting power
• NZIF receiver with –94 dBm Bluetooth LE sensitivity
• Adaptive Frequency Hopping (AFH)
• Standard HCI based on SDIO/SPI/UART
• High-speed UART HCI, up to 4 Mbps
• Bluetooth 4.2 BR/EDR and Bluetooth LE dual mode controller
• Synchronous Connection-Oriented/Extended (SCO/eSCO)
• CVSD and SBC for audio codec
• Bluetooth Piconet and Scatternet
• Multi-connections in Classic Bluetooth and Bluetooth LE
• Simultaneous advertising and scanning
CPU and Memory
• Xtensa® single-/dual-core 32-bit LX6 microprocessor(s)
• CoreMark® score:
– 1 core at 240 MHz: 539.98 CoreMark; 2.25 CoreMark/MHz
Espressif Systems 3
Submit Documentation Feedback
ESP32 Series Datasheet v5.2
– 2 cores at 240 MHz: 1079.96 CoreMark; 4.50 CoreMark/MHz
• 448 KB ROM
• 520 KB SRAM
• 16 KB SRAM in RTC
• QSPI supports multiple flash/SRAM chips
Clocks and Timers
• Internal 8 MHz oscillator with calibration
• Internal RC oscillator with calibration
• External 2 MHz ~ 60 MHz crystal oscillator (40 MHz only for Wi-Fi/Bluetooth functionality)
• External 32 kHz crystal oscillator for RTC with calibration
• Two timer groups, including 2 × 64-bit timers and 1 × main watchdog in each group
• One RTC timer
• RTC watchdog
Advanced Peripheral Interfaces
• 34 programmable GPIOs
– Five strapping GPIOs
– Six input-only GPIOs
– Six GPIOs needed for in-package flash (ESP32-U4WDH) and in-package PSRAM
(ESP32-D0WDRH2-V3)
• 12-bit SAR ADC up to 18 channels
• Two 8-bit DAC
• 10 touch sensors
• Four SPI interfaces
• Two I2S interfaces
• Two I2C interfaces
• Three UART interfaces
• One host (SD/eMMC/SDIO)
• One slave (SDIO/SPI)
• Pulse count controller
• Ethernet MAC interface with dedicated DMA and IEEE 1588 support
• TWAI®, compatible with ISO 11898-1 (CAN Specification 2.0)
• RMT (TX/RX)
Espressif Systems 4
Submit Documentation Feedback
ESP32 Series Datasheet v5.2
• Motor PWM
• LED PWM up to 16 channels
Power Management
• Fine-resolution power control through a selection of clock frequency, duty cycle, Wi-Fi operating modes,
and individual power control of internal components
• Five power modes designed for typical scenarios: Active, Modem-sleep, Light-sleep, Deep-sleep,
Hibernation
• Power consumption in Deep-sleep mode is 10 μA
• Ultra-Low-Power (ULP) coprocessors
• RTC memory remains powered on in Deep-sleep mode
Security
• Secure boot
• Flash encryption
• 1024-bit OTP, up to 768-bit for customers
• Cryptographic hardware acceleration:
– AES
– Hash (SHA-2)
– RSA
– Random Number Generator (RNG

Applications
With low power consumption, ESP32 is an ideal choice for IoT devices in the following areas:
• Smart Home
• Industrial Automation
• Health Care
• Consumer Electronics
• Smart Agriculture
• POS Machines
• Service Robot
• Audio Devices
• Generic Low-power IoT Sensor Hubs
• Generic Low-power IoT Data Loggers
• Cameras for Video Streaming
• Speech Recognition
• Image Recognition
• SDIO Wi-Fi + Bluetooth Networking C 




